%!TEX root = ../template.tex
%%%%%%%%%%%%%%%%%%%%%%%%%%%%%%%%%%%%%%%%%%%%%%%%%%%%%%%%%%%%%%%%%%%%
%% abstrac-pt.tex
%% NOVA thesis document file
%%
%% Resumo em Português
%%%%%%%%%%%%%%%%%%%%%%%%%%%%%%%%%%%%%%%%%%%%%%%%%%%%%%%%%%%%%%%%%%%%

\typeout{NT FILE abstrac-pt.tex}%

Os serviços \emph{online} modernos servem centenas de milhões de
utilizadores a partir de centros de dados distribuídos por todo o mundo,
tendo de conciliar três requisitos em tensão: baixa latência, elevada
disponibilidade e consistência dos dados. À medida que o número de centros
de dados aumenta, replicar todos os objetos em todos os centros torna-se
proibitivo em termos de custos de armazenamento, rede e processamento. A
\emph{replicação parcial}, em que cada centro de dados replica apenas um
subconjunto dos dados, é uma forma natural de conter estes custos, mas
introduz desafios significativos: manter invariantes da aplicação sob
consistência fraca, executar transações cujos objetos não estão todos
replicados localmente e responder a consultas analíticas recorrentes
que dependem de dados dispersos por vários centros de dados.

Esta tese investiga como conceber uma base de dados geo-distribuída que
suporte replicação parcial de forma eficiente, oferecendo simultaneamente
garantias de correção úteis e uma interface rica e expressiva para as
aplicações. Propomos a PotionDB, uma base de dados geo-distribuída cujas
contribuições assentam em três pilares: (i) um algoritmo de replicação
que combina replicação parcial, não-uniforme e geo-replicação, propagando
para cada servidor apenas as operações relevantes para os dados que
replica localmente; (ii) um algoritmo transacional que oferece consistência
PSI+cset, com \emph{snapshot isolation} dentro de cada servidor e
consistência causal entre servidores, suportando transações cujo estado
pode não estar totalmente presente numa única réplica; e (iii) vistas
materializadas mantidas incrementalmente através de operações sobre CRDTs,
mantidas atomicamente consistentes com os objetos que sumarizam, capazes
de resumir dados particionados por várias réplicas parciais. A PotionDB
inclui ainda um algoritmo de reconstrução de versões antigas dos objetos
baseado em registos de operações e efeitos, reduzindo o consumo de memória
relativamente à manutenção do histórico completo de versões.

O trabalho futuro foca-se em particionamento dinâmico, níveis de
consistência por transação (incluindo invariantes sobre objetos e vistas),
uma interface do tipo SQL e um algoritmo refinado de \emph{garbage
collection}. O sistema proposto está implementado num protótipo
funcional e é avaliado no Grid'5000 com recurso ao TPC-H, ao YCSB e a
\emph{microbenchmarks}, demonstrando que consultas analíticas complexas
sobre dados parcialmente replicados podem ser respondidas com latência
comparável à de leituras simples de objetos, preservando a consistência
entre as vistas e os objetos a que se referem.

% Palavras-chave do resumo em Português
\begin{keywords}
Bases de dados geo-distribuídas, Replicação parcial, Replicação
não-uniforme, Vistas materializadas, CRDTs, Consistência causal,
\emph{Snapshot isolation}, Transações, Invariantes
\end{keywords}
